\chapter{Metodología II: arquitecturas, entrenamiento y evaluación}
\label{cap:metodologia-modelos}
\section{Introducción}
Se describieron las arquitecturas y sus representaciones de entrada, que determinaron los modelos comparados; el esquema de entrenamiento fijó sus configuraciones y pesos. Los protocolos de evaluación produjeron exactitudes por sujeto en escenarios intra-sujeto, inter-sujeto mediante dejar un sujeto fuera (Leave-One-Subject-Out, \loso{}) y de transferencia. Los indicadores y el análisis estadístico definieron las medidas de degradación y comparación; los criterios adicionales establecieron medidas de costo, estabilidad, robustez, interpretabilidad y eficiencia de datos para complementar la evaluación.
\section{Fase 3: arquitecturas e implementación}
\label{sec:fase3}

Se implementaron y evaluaron cuatro arquitecturas para señales de electroencefalografía (EEG): la de
referencia (\eegnet{}), el Conformer, el modelo fundacional LaBraM (variante Base) y
EEGMamba, cuya comparación se presenta en la sección~\ref{sec:comparacion}. EEGMamba
requirió las adaptaciones específicas descritas en la
sección~\ref{sec:eegmamba-adapt}.

LaBraM Base fue el único modelo fundacional empleado: se utilizaron pesos preentrenados y ajuste fino completo. EEGMamba se trató como una arquitectura de espacio de estados y se entrenó desde cero, al igual que el Conformer.

\subsection{Representación de entrada y procesamiento por arquitectura}
\label{sec:representaciones}

Las cuatro arquitecturas reciben la misma señal preprocesada, pero no la consumen de la
misma forma. Esta sección detalla qué hace cada una con ella, porque la
interpretación de los resultados depende de ello: dos de las cuatro operan sobre la forma
de onda continua y las otras dos la fragmentan en parches de un segundo. Las dimensiones que
se indican, $62$ canales por $801$ muestras, son las de los protocolos sobre \openbmi{}; en
el capítulo~\ref{cap:validacion-externa} la entrada tiene los $45$ canales comunes a los
dos conjuntos, y el número de parámetros de la capa que depende de los canales cambia en
consecuencia.

\subsubsection{\eegnet{}: convolución separable sobre la señal continua}

\eegnet{} \cite{lawhern2018}, en la implementación de \texttt{braindecode} 1.6.1
\cite{schirrmeister2017} con sus hiperparámetros por defecto, consume el tensor $(C \times T) = (62 \times 801)$ sin ninguna transformación intermedia.
Su primer bloque aplica convoluciones \emph{temporales}, cuyos núcleos actúan como filtros
de banda aprendidos; el segundo aplica una convolución \emph{en profundidad} sobre el eje
de canales, un filtro espacial por cada filtro temporal, que cumple un papel análogo al de
los patrones espaciales comunes; el tercero combina ambos ejes con una convolución
separable. Un promediado temporal y una capa lineal producen los logits
(figura~\ref{fig:arq-eegnet}). Con $62$ canales
tiene $2\,898$ parámetros, de dos a algo más de tres órdenes de magnitud menos que las demás.

\begin{figure}[H]
\centering
\begin{tikzpicture}[
  font=\scriptsize,
  b/.style={draw, rounded corners=2pt, minimum height=0.7cm, align=center, fill=blue!8},
  fl/.style={-{Latex[length=1.6mm]}, semithick}, x=1cm, y=1cm
]
\node[b, minimum width=1.9cm] (in) at (0,0) {$62 \times 801$\\señal continua};
\node[b, minimum width=2.0cm] (t)  at (2.6,0) {conv.\ temporal\\(filtros de banda)};
\node[b, minimum width=2.0cm] (e)  at (5.3,0) {conv.\ en profundidad\\(filtros espaciales)};
\node[b, minimum width=1.9cm] (sp) at (7.9,0) {conv.\ separable\\+ promediado};
\node[b, minimum width=1.2cm, fill=orange!18] (o) at (10.0,0) {lineal\\2 clases};
\foreach \a/\bb in {in/t,t/e,e/sp,sp/o} \draw[fl] (\a) -- (\bb);
\end{tikzpicture}
\caption{\eegnet{}. Opera sobre la señal continua; separa el filtrado temporal del
espacial, en ese orden.}
\label{fig:arq-eegnet}
\end{figure}

\subsubsection{Conformer: convolución seguida de atención}

El EEG-Conformer \cite{song2023conformer}, también en la implementación de
\texttt{braindecode} con sus hiperparámetros por defecto, recibe el mismo tensor
$(62 \times 801)$ que \eegnet{} y no necesita adaptador. Un bloque
convolucional inicial, temporal y después espacial, transforma la señal en una secuencia de
vectores o \emph{tokens}; sobre esa secuencia actúa un codificador con autoatención, que
pondera cada instante en función de todos los demás en lugar de hacerlo dentro de una
ventana local fija como la convolución. Un perceptrón produce los logits. Tiene $710\,146$
parámetros (figura~\ref{fig:arq-conformer}).

\begin{figure}[H]
\centering
\begin{tikzpicture}[
  font=\scriptsize,
  b/.style={draw, rounded corners=2pt, minimum height=0.7cm, align=center, fill=blue!8},
  fl/.style={-{Latex[length=1.6mm]}, semithick}, x=1cm, y=1cm
]
\node[b, minimum width=1.9cm] (in) at (0,0) {$62 \times 801$\\señal continua};
\node[b, minimum width=2.4cm] (c)  at (2.9,0) {conv.\ temporal\\+ conv.\ espacial};
\node[b, minimum width=1.9cm, fill=green!12] (tk) at (5.7,0) {secuencia de\\\emph{tokens}};
\node[b, minimum width=2.1cm] (at) at (8.2,0) {autoatención\\(global)};
\node[b, minimum width=1.2cm, fill=orange!18] (o) at (10.4,0) {MLP\\2 clases};
\foreach \a/\bb in {in/c,c/tk,tk/at,at/o} \draw[fl] (\a) -- (\bb);
\end{tikzpicture}
\caption{Conformer. La convolución construye los \emph{tokens} y la autoatención los
relaciona entre sí sin restricción de distancia temporal.}
\label{fig:arq-conformer}
\end{figure}

\subsubsection{LaBraM: parches de un segundo con identidad de canal}

LaBraM \cite{jiang2024labram} es el único modelo con pesos preentrenados. Divide la ventana
en parches de $200$ muestras,
es decir $1$\,s a \hz{200}, lo que obliga a truncar de $801$ a $800$ y produce $4$ parches
por canal. Cada parche se proyecta a un vector, al que se suma un \textbf{embedding de la
identidad del canal}: la entrada incluye explícitamente el electrodo de origen
de cada parche, información que no forma parte de la entrada de las tres arquitecturas restantes. Esa lista
de índices se toma de un vocabulario de $128$ posiciones de electrodo, que incluye
posiciones de los sistemas 10-10 y 10-05, y es lo que permite aplicar el
modelo a montajes distintos. Un codificador \emph{transformer} procesa el conjunto de
parches. Se emplea la variante Base, $5\,818\,138$ parámetros, con ajuste fino completo,
partiendo de los pesos distribuidos en \texttt{braindecode/labram-pretrained}
(figura~\ref{fig:arq-labram}).

Tres decisiones de esta adaptación se apartan del procedimiento con que se preentrenó el
modelo y condicionan su evaluación. La señal se le entrega normalizada por canal
(\textit{z-score}), igual que a las demás arquitecturas, y no en microvoltios divididos por
$100$ como en el preentrenamiento. El \emph{embedding} temporal preentrenado, de $16$
posiciones, se interpola a las $5$ que exige la ventana de $4$\,s. La capa de clasificación
se reinicia, y el ajuste fino usa una tasa de aprendizaje uniforme de $10^{-4}$, sin el
decaimiento por capa del procedimiento original.

\begin{figure}[H]
\centering
\begin{tikzpicture}[
  font=\scriptsize,
  b/.style={draw, rounded corners=2pt, minimum height=0.7cm, align=center, fill=blue!8},
  fl/.style={-{Latex[length=1.6mm]}, semithick}, x=1cm, y=1cm
]
\node[b, minimum width=1.9cm] (in) at (0,0) {$62 \times 801$\\señal continua};
\node[b, minimum width=2.2cm, fill=green!12] (pa) at (2.9,0)
  {parches de \seg{1}\\$62 \times 4 \times 200$};
\node[b, minimum width=2.0cm] (em) at (5.7,0) {proyección\\de cada parche};
\node[b, minimum width=2.0cm] (tr) at (8.3,0) {\emph{transformer}\\preentrenado};
\node[b, minimum width=1.2cm, fill=orange!18] (o) at (10.5,0) {lineal\\2 clases};
\node[b, minimum width=2.4cm, fill=red!10] (ch) at (5.7,-1.4)
  {índices de canal\\(vocabulario de 128)};
\foreach \a/\bb in {in/pa,pa/em,em/tr,tr/o} \draw[fl] (\a) -- (\bb);
\draw[fl] (ch) -- (em);
\end{tikzpicture}
\caption{LaBraM. Único modelo cuya entrada incluye la identidad del electrodo junto a la señal, lo
que le permite operar sobre montajes distintos del de preentrenamiento.}
\label{fig:arq-labram}
\end{figure}

\subsubsection{EEGMamba: espacio de estados sobre parches, con rama espectral}

Se emplea la implementación pública de EEGMamba \cite{wang2025eegmamba} como arquitectura de espacio de estados entrenada desde cero. Fragmenta la ventana en los mismos $4$
parches de $200$ muestras. Cada parche se codifica
por dos vías que se suman: una \emph{temporal}, una convolución de $25$ filtros de
longitud $49$ y paso $25$, y una
\emph{espectral}, que toma el módulo de la transformada rápida de Fourier real del parche,
$101$ coeficientes, y lo proyecta al mismo espacio; se añade una codificación posicional
convolucional. Es la única de las cuatro cuya entrada incluye
una representación en frecuencia de forma explícita. La secuencia resultante se procesa
mediante una pila de $12$ bloques Mamba2 \cite{gu2023mamba,dao2024mamba2}, modelos de
espacio de estados aplicados en ambos sentidos de la secuencia, cuyo
costo crece de forma lineal con la longitud de la secuencia y no cuadrática como el de la
autoatención. Se emplea la cabeza que conserva la representación por parche y canal;
tiene $3\,376\,445$ parámetros y se entrena desde cero
(figura~\ref{fig:arq-eegmamba}).

\begin{figure}[H]
\centering
\begin{tikzpicture}[
  font=\scriptsize,
  b/.style={draw, rounded corners=2pt, minimum height=0.7cm, align=center, fill=blue!8},
  fl/.style={-{Latex[length=1.6mm]}, semithick}, x=1cm, y=1cm
]
\node[b, minimum width=1.7cm] (in) at (0,0) {$62 \times 801$\\señal continua};
\node[b, minimum width=2.0cm, fill=green!12] (pa) at (2.6,0)
  {parches de \seg{1}\\$62 \times 4 \times 200$};
\node[b, minimum width=1.7cm] (tt) at (5.1,0.75) {proyección\\temporal};
\node[b, minimum width=1.7cm, fill=red!10] (ff) at (5.1,-0.75) {$|\text{rFFT}|$\\(101 coef.)};
\node[b, minimum width=2.2cm] (mm) at (7.9,0) {Mamba2\\bidireccional};
\node[b, minimum width=1.2cm, fill=orange!18] (o) at (10.2,0) {lineal\\2 clases};
\draw[fl] (in) -- (pa);
\draw[fl] (pa) |- (tt);
\draw[fl] (pa) |- (ff);
\draw[fl] (tt) -| (mm);
\draw[fl] (ff) -| (mm);
\draw[fl] (mm) -- (o);
\node[font=\scriptsize] at (6.6,0) {$+$};
\end{tikzpicture}
\caption{EEGMamba. Suma una codificación temporal y una espectral de cada parche, y procesa
la secuencia con bloques de espacio de estados en ambos sentidos.}
\label{fig:arq-eegmamba}
\end{figure}

Resumiendo el contraste que importa para la comparación: \eegnet{} y el Conformer operan
sobre la señal continua, mientras que LaBraM y EEGMamba la fragmentan en parches de un
segundo. LaBraM es el único cuya entrada incluye la identidad del electrodo y el único preentrenado;
EEGMamba es el único cuya entrada incluye una representación en frecuencia explícita. Ninguna de las
cuatro abandona el dominio del tiempo como espacio de predicción, punto sobre el que vuelve
la discusión.

\subsection{Esquema de entrenamiento unificado}

Todas las arquitecturas comparten el mismo esquema, con una excepción que se declara
más abajo (el recorte de la norma del gradiente), para que la comparación
atribuya las diferencias al modelo y no al ajuste: optimizador AdamW
\cite{loshchilov2019adamw} con decaimiento de pesos de $0.01$,
calentamiento lineal de $10$ épocas seguido de decaimiento coseno calculado sobre el
máximo de $200$ épocas, suavizado de etiquetas de $0.1$ en la pérdida de entrenamiento
(la de validación se calcula sin suavizado),
precisión mixta automática, tamaño de lote efectivo fijo \emph{dentro de cada
arquitectura}, y parada temprana por pérdida de validación con paciencia de 20 épocas
sobre un máximo de 200. Al terminar se restauran los pesos de la época de menor pérdida
de validación.

Dos consecuencias de esta configuración no son evidentes y afectan a la lectura de los
resultados. La paciencia no corre durante el calentamiento, de modo que la parada más
temprana posible es la época $30$. Y como el coseno se calcula sobre $200$ épocas mientras
las ejecuciones \loso{} se detienen entre las épocas $30$ y $146$, la tasa de aprendizaje
decae poco antes de la parada en la mayoría de ellas: conserva el $97$\,\% de su valor en la
época $30$ y el $19$\,\% en la $146$. El cómputo no es determinista bit a bit (cuDNN sin modo determinista y precisión
mixta), así que repetir una ejecución con la misma semilla puede variar en los últimos
decimales. Las versiones empleadas son \texttt{braindecode} 1.6.1, MNE-Python 1.12.1,
\texttt{mne-icalabel} 0.9.0 y MOABB 1.5.0 \cite{jayaram2018moabb}. Las ejecuciones se
agrupan en dos regímenes de entrenamiento. El \emph{primer régimen} reúne las ejecuciones
previas (unidad de procesamiento gráfico (GPU) NVIDIA A40; recorte de la norma del gradiente solo en EEGMamba); el
\emph{segundo régimen}, las posteriores (GPU NVIDIA L40S; recorte a $1.0$ en las cuatro
arquitecturas). La
tabla~\ref{tab:regimen} indica el régimen de cada semilla y la tabla~\ref{tab:hardware} el
hardware; esta denominación se usa en el resto del documento.

\textbf{Recorte de la norma del gradiente.}\label{sec:recorte-gradiente}
En el primer régimen, las ejecuciones base intra-sujeto y \loso{} de \eegnet{}, el
Conformer y LaBraM se entrenaron \emph{sin} recorte de la norma del gradiente, y las de
EEGMamba con recorte a $1.0$. En el segundo régimen las cuatro arquitecturas usan recorte a
$1.0$. La transferencia de la semilla $2026$ también lo llevó, aunque se entrenó en A40
(tabla~\ref{tab:regimen}). Cada ejecución guarda un identificador de configuración
(\textit{hash}) que permite comprobar si el recorte estuvo activo.

El recorte se usa en el segundo régimen porque los núcleos de espacio de estados de EEGMamba
saturan el rango de \texttt{float16} y divergen sin él, y es una técnica descrita para los
gradientes explosivos \cite{pascanu2013difficulty}; aplicarlo a las cuatro arquitecturas
mantiene un único procedimiento de entrenamiento. No se comprobó que fuese neutro para
\eegnet{}, el Conformer y LaBraM, y no se conserva registro numérico de la divergencia de
EEGMamba sin recorte ni de la norma del gradiente, por lo que se desconoce en qué fracción de
los pasos el recorte actúa. Las recetas de referencia de las cuatro arquitecturas no
coinciden en usarlo: el repositorio de EEGMamba \cite{wang2025eegmamba} lo activa por defecto
con valor $1$; el de LaBraM \cite{jiang2024labram} lo deja desactivado en el ajuste fino; y
en los de \eegnet{} \cite{lawhern2018} y del Conformer \cite{song2023conformer} no aparece.
Este trabajo no afirma que el recorte sea una práctica habitual en el decodificado de EEG.

En \eegnet{}, el Conformer y LaBraM, el contraste entre una variante (segundo régimen) y la
ejecución base de la semilla $2026$ mezcla el efecto de la variante con un cambio compuesto
del régimen de entrenamiento: recorte de gradiente, hardware (A40 frente a L40S) y versión
del código. En EEGMamba el recorte no cambia, aunque sí el hardware y el código. Ese cambio
no puede aislarse en sus componentes, pero se acota su magnitud comparando, para cada sujeto,
la semilla $2026$ de la base con la media de las semillas $2027$ y $2028$
(sección~\ref{sec:cota-confusor}; resultados en la sección~\ref{sec:confusor}). No se detecta
efecto del cambio de régimen en ninguna arquitectura: el límite del intervalo de confianza al
$95$\,\% de la diferencia media por sujeto en exactitud es de $0.010$ a $0.020$ según la
arquitectura, y en una media de tres semillas, de las cuales una se entrenó sin recorte, la
contaminación máxima es de $0.003$ a $0.007$ (tabla~\ref{tab:confusor}). Eso descarta un
efecto medio mayor que esos límites, no un efecto nulo. Para el protocolo intra-sujeto de
\eegnet{} y del Conformer no existe base con recorte (las tres semillas se entrenaron sin él),
de modo que ese efecto no es estimable y se declara en las limitaciones. Repetir esa base con
recorte queda como trabajo pendiente.

\textbf{Alcance de la unificación.} Las cuatro arquitecturas comparten la \emph{estructura}
del procedimiento (optimizador, planificación de la tasa de aprendizaje, criterio de parada,
aumentación, particiones y métricas) y no los valores numéricos, que enumera la
tabla~\ref{tab:configs}.

\begin{table}[H]
\centering
\caption{Configuración efectiva de cada arquitectura. Las dos primeras columnas son los
valores con los que se ejecutaron los resultados reportados; la última resume por qué
difieren del resto. La tabla no recoge el recorte de la norma del gradiente, que no es un
valor por arquitectura sino un cambio de régimen entre ejecuciones (párrafo
``Recorte de la norma del gradiente'' de esta sección).}
\label{tab:configs}
\small
\begin{tabular}{lccp{7.2cm}}
\toprule
 & \textbf{Tasa de aprendizaje} & \textbf{Lote} & \textbf{Motivo de la desviación} \\
\midrule
\eegnet{}   & $2.8\times10^{-3}$ & $256$ & Empleada al pasar al lote $256$, sin búsqueda ni criterio de ajuste registrado; coincide con escalar $10^{-3}$ por la raíz cuadrada del aumento de lote, $10^{-3}\sqrt{8}$ \\
Conformer   & $1.0\times10^{-3}$ & $256$ & Valor por defecto del esquema; no se comparó de forma sistemática con el de \eegnet{} \\
LaBraM      & $1.0\times10^{-4}$ & $32$  & Parte de pesos preentrenados: en una prueba preliminar, con la tasa por defecto producía una salida constante para todas las entradas \\
EEGMamba    & $1.0\times10^{-3}$ & $32$  & Cómputo en \texttt{bfloat16}: los núcleos de espacio de estados superan el rango numérico representable en \texttt{float16} y divergen (sección~\ref{sec:eegmamba-adapt}) \\
\bottomrule
\end{tabular}
\end{table}

Los ajustes responden a problemas de entrenamiento observados en pruebas preliminares y no
a la búsqueda de un mejor resultado para una arquitectura concreta; no se realizó una
búsqueda sistemática de hiperparámetros. Sí hubo decisiones tomadas observando el conjunto
principal: la cabeza de EEGMamba se eligió comparando su exactitud intra-sujeto sobre las
mismas particiones que luego se reportan (sección~\ref{sec:eegmamba-adapt}) y se trasladó al
protocolo \loso{}, lo que pudo favorecer ligeramente a esa arquitectura, y la tasa de LaBraM
y la precisión \texttt{bfloat16} se fijaron tras pruebas preliminares sobre el mismo
conjunto. La comparación debe leerse, por tanto, entre arquitecturas en una configuración
funcional y no óptima. El tamaño de lote tampoco es constante ($256$ frente a $32$), lo que
altera el número de actualizaciones del optimizador por época y no solo la velocidad de
cómputo.

Se tomaron como referencia los valores de los artículos originales o de sus implementaciones, con las excepciones documentadas en la tabla~\ref{tab:configs} y las adaptaciones descritas en esta sección, incluidas las tasas de \eegnet{} y EEGMamba y la cabeza de EEGMamba. La tasa de \eegnet{} carece de un criterio de selección registrado, como indica la tabla. El ajuste fino de pesos preentrenados solo se aplicó a LaBraM. Una sintonización específica por arquitectura podría alterar su orden; esta limitación se aborda en la sección~\ref{sec:limitaciones}.

\textbf{Aumentación de datos.} Todos los entrenamientos, en las cuatro arquitecturas y
en los dos protocolos, aplican tres transformaciones sobre el conjunto de entrenamiento,
nunca sobre validación ni sobre prueba. Se aplican \emph{por minilote}, de modo que cada
época ve realizaciones distintas:

\begin{itemize}
  \item \textbf{Desplazamiento temporal} de hasta el $5$\,\% de la ventana, es decir
  $\pm 40$ muestras o $\pm 200$\,ms, sorteado de forma independiente por segmento.
  \item \textbf{Ruido gaussiano} de desviación igual al $2$\,\% de la desviación típica
  \emph{de cada canal}, calculada sobre el propio lote.
  \item \textbf{Enmascarado en frecuencia} de una banda aleatoria de hasta $2$\,Hz,
  aplicada sobre la transformada de Fourier y deshecha con la transformada inversa.
\end{itemize}

La mezcla de segmentos (\textit{mixup}) está implementada pero desactivada. La
aumentación es idéntica en las cuatro arquitecturas y es la línea base frente a la que se
contrastan las variantes de aumentación. Es \textbf{en línea}: el conjunto de entrenamiento
conserva su tamaño (por ejemplo, los $18\,360$ segmentos de cada \textit{fold} \loso{}) y
solo varía entre épocas la perturbación con que cada segmento se presenta a la red, de modo
que actúa como regularización y no amplía el número de ejemplos. Solo una de las tres
transformaciones añade ruido; el desplazamiento temporal es geométrico y el enmascarado en
frecuencia suprime una banda.

La arquitectura de referencia es \eegnet{} \cite{lawhern2018}, el punto de comparación
habitual en la literatura de imaginación motora.

\subsection{Adaptaciones específicas de EEGMamba}
\label{sec:eegmamba-adapt}

A diferencia de LaBraM, EEGMamba se entrenó desde cero, sin pesos preentrenados.
Por ello se aparta del esquema unificado en tres puntos, todos necesarios para que
el modelo aprenda y ninguno orientado a favorecerlo frente al resto de
arquitecturas:

\begin{itemize}
  \item \textbf{Precisión de cómputo.} Los núcleos de espacio de estados del
  bloque Mamba (recorrido selectivo y exponenciales) superan el rango numérico
  representable en media precisión \texttt{float16} y producen valores no numéricos pasadas las primeras
  épocas. Se sustituye \texttt{float16} por \texttt{bfloat16}, que conserva el
  rango de exponente de la precisión simple y elimina el desbordamiento sin
  penalización de velocidad apreciable.
  \item \textbf{Recorte de gradiente.} La norma del gradiente se acota a $1.0$
  \cite{pascanu2013difficulty}. En las ejecuciones base del primer régimen es una
  adaptación propia de esta arquitectura; en el segundo régimen lo llevan las cuatro. Su
  justificación y su consecuencia para la comparación se tratan en el párrafo ``Recorte de
  la norma del gradiente'' de la subsección anterior.
  \item \textbf{Cabeza de clasificación.} Se emplea en \textbf{ambos protocolos} la
  cabeza que conserva la representación de todos los canales y parches y la proyecta
  linealmente, porque es la que preserva la lateralización sensoriomotora, el contraste
  que separa la imaginación motora de una mano y de la otra. La alternativa considerada,
  un promediado global sobre canales y parches, elimina la información de lateralización antes de la
  proyección, y con ella la exactitud intra-sujeto queda en $0.519$, apenas por encima
  del azar (Wilcoxon de una muestra frente a $0.50$, $p = 0.017$). Conservando la
  representación por canal asciende a $0.577$ con la semilla 2026 ($p = 1.5\times10^{-8}$ frente al azar) y
  supera a la variante de promediado con tamaño de efecto grande
  ($p = 2.6\times10^{-6}$, $r = 0.64$; contraste de Wilcoxon pareado sobre los 54
  sujetos).

  Esta cabeza sobreajusta en el protocolo intra-sujeto. Con 62 canales la proyección
  recibe $62 \times 4 \times 200 = 49\,600$ características frente a los 144 segmentos de
  entrenamiento de cada partición, y los registros por época muestran que la pérdida de
  validación aumenta desde las primeras épocas mientras la de entrenamiento disminuye.
  La consecuencia es que, en el $85.6$\,\% de las ejecuciones intra-sujeto, el criterio de
  parada conserva un modelo de las épocas de calentamiento
  (sección~\ref{sec:seleccion-intra}). La exactitud de $0.577$ debe leerse, por tanto,
  como el resultado conjunto de la arquitectura y del procedimiento de selección del
  modelo con 36 ensayos de validación, y no como una medida de capacidad comparable con
  la de las otras tres arquitecturas.
\end{itemize}

La tasa de aprendizaje se fija en $10^{-3}$, superior a la empleada al ajustar el
modelo preentrenado, coherente con un entrenamiento desde cero.

\subsection{Semillas y régimen de optimización de cada una}
\label{sec:semillas-metodo}\label{sec:fase3-variantes}

Las cuatro arquitecturas se ejecutaron con tres semillas, $2026$, $2027$ y $2028$, en el
protocolo intra-sujeto y en \loso{}, y las tres semillas se repitieron también en el modelo
de $54$ sujetos transferido a WBCIC-SHU, en el modelo nativo de WBCIC-SHU y en el protocolo
de sujetos apartados (44 más 10). Eso significa, por protocolo:

\begin{itemize}
  \item \textbf{Intra-sujeto.} \eegnet{} y el Conformer ya contaban con las tres semillas;
  para LaBraM y EEGMamba se añadieron la $2027$ y la $2028$. Cada arquitectura aporta
  $54 \times 20 \times 3 = 3\,240$ mediciones.
  \item \textbf{\loso{}.} La semilla $2026$ procede de las ejecuciones del primer régimen, y se
  añadieron la $2027$ y la $2028$ para las cuatro arquitecturas: $54 \times 3 = 162$
  \textit{folds} por arquitectura.
  \item \textbf{Transferencia y sujetos apartados.} Se añadieron las semillas $2027$ y
  $2028$ a las tres configuraciones que antes tenían la $2026$.
\end{itemize}

El número de semillas se limita a tres por el presupuesto de cómputo. Se reportan todas; no
se elige la de mejor desempeño, porque seleccionarla sobre el propio conjunto de prueba
sería un sesgo de selección. Las variantes de preprocesamiento y de aumentación
(secciones~\ref{sec:variantes-preproc} y~\ref{sec:variantes-aug}) usan solo la semilla
$2026$, para cubrir más variantes con una semilla cada una; por ello el desplazamiento medio
entre semillas no se incorpora en esos contrastes (sección~\ref{sec:ablaciones-f4}).

Las semillas controlan la inicialización de la red y los sorteos del entrenamiento. Se
emparejan por índice entre arquitecturas, aunque las inicializaciones son independientes, de
modo que el emparejamiento solo etiqueta. La semilla del remuestreo estadístico es otra
(sección~\ref{sec:bootstrap}).



\begin{table}[H]
\centering
\caption{Régimen de recorte de gradiente por semilla (orden $2026$, $2027$, $2028$).
S: sin recorte. C: con recorte a $1.0$. En las semillas $2026$ de las cuatro arquitecturas,
las ejecuciones son del primer régimen (GPU A40); las demás, del segundo (GPU
L40S), salvo las semillas $2027$ y $2028$ intra-sujeto de \eegnet{} y del Conformer, que
también se entrenaron en A40. La transferencia de la semilla $2026$, aunque anterior al segundo régimen y
en A40, ya llevaba recorte.}
\label{tab:regimen}
\small
\begin{tabular}{lcccc}
\toprule
 & \textbf{\eegnet{}} & \textbf{Conformer} & \textbf{LaBraM} & \textbf{EEGMamba} \\
\midrule
\loso{}                          & S C C & S C C & S C C & C C C \\
Intra-sujeto                     & S S S & S S S & S C C & C C C \\
Transferencia y 44 + 10          & C C C & C C C & C C C & C C C \\
\bottomrule
\end{tabular}
\end{table}

La tabla muestra qué comparaciones cruzan regímenes: las de \eegnet{}, el Conformer y
LaBraM entre la semilla $2026$ y las otras dos en \loso{}, y las de LaBraM en
intra-sujeto. Entre las semillas de una misma fila sin cambio de letra no hay cambio de
recorte, aunque sí puede haber cambio de hardware.



\subsection{Variantes de aumentación de datos}
\label{sec:variantes-aug}

La configuración base aplica las tres transformaciones descritas en la
sección~\ref{sec:fase3} (desplazamiento temporal, ruido gaussiano y enmascarado en
frecuencia). Se evalúan tres cambios, todos sobre el conjunto de entrenamiento y por
minilote, como la base:

\begin{itemize}
  \item \textbf{Sin aumentación.} Las tres transformaciones de la base en cero. Esta
  condición no es la base: es la ausencia de toda aumentación.
  \item \textbf{Reflexión de canales} \cite{wang2024channelreflection}. La base más una
  reflexión que, con probabilidad $0.5$ por segmento, intercambia cada canal con su simétrico
  del otro hemisferio y permuta la etiqueta (mano izquierda por mano derecha), porque la
  desincronización de la imaginación motora es contralateral. El simétrico de un canal impar
  $n$ es el par $n+1$, según el sistema 10-20 y sus extensiones; los canales de la línea media
  y los que no tienen pareja en el montaje se quedan en su sitio. Se aplica antes que el
  resto de transformaciones del lote.
  \item \textbf{Ruido en la amplitud de Fourier.} La base más una perturbación
  multiplicativa gaussiana del módulo de cada componente de la transformada de Fourier,
  $|F'(f)| = |F(f)|\,\max\!\bigl(0,\ 1 + \sigma\,\mathcal{N}(0,1)\bigr)$ con $\sigma = 0.1$,
  independiente por canal, segmento y componente, con la fase conservada, y se vuelve al
  tiempo con la transformada inversa. No se dispuso de un valor de referencia en la
  literatura para esta aplicación; $\sigma = 0.1$ es una elección propia, fijada antes de
  ejecutar y sin ajuste posterior, por lo que el resultado vale para esa amplitud y no para
  el ruido en Fourier en general.
\end{itemize}

En el protocolo intra-sujeto se ejecutan las tres con la semilla $2026$; en \loso{}, solo la
reflexión de canales, con la semilla $2026$. El ruido en Fourier en \loso{} no se ejecutó.

\subsection{Variantes con una sola semilla y ablaciones recortadas en LaBraM y EEGMamba}
\label{sec:ablaciones-recortadas}

LaBraM y EEGMamba cuestan varias horas cada ablación intra-sujeto, porque
saturan la GPU y la paralelización por trabajadores apenas las acelera. Por coste de
cómputo se decidió, antes de ejecutarlas, limitar sus ablaciones a la reflexión de canales,
el ruido en Fourier y la ventana de \seg{0} a \seg{2}, esta última porque cambia el número
de parches. \eegnet{} y el Conformer recorren las ocho: tres subbandas, dos ventanas, sin
aumentación, reflexión de canales y ruido en Fourier. Por tanto los efectos de subbanda,
de la ventana de \seg{3} y de la ausencia de aumentación no se evalúan en LaBraM ni en
EEGMamba.



\section{Fase 4: protocolo de evaluación}
\label{sec:fase4}

Se emplean dos protocolos de validación, que responden a preguntas distintas. El
protocolo intra-sujeto mide cuánta información motora puede decodificarse en un
sujeto a partir de sus propios datos de entrenamiento. El protocolo inter-sujeto
mide cuánto de ese aprendizaje se transfiere a un sujeto nuevo. La comparación
entre ambos es la que motiva el indicador de la sección~\ref{sec:dsi}.

\subsection{Validación intra-sujeto}
\label{sec:within}

Cada sujeto se evaluó consigo mismo y \textbf{por separado en cada sesión}, para
que el resultado no mezcle la variabilidad entre días con la del propio
protocolo. Los 200 ensayos de una sesión se reparten en diez \textit{folds} estratificados por
clase, sorteados con barajado dentro de cada sesión (\texttt{StratifiedKFold}), y no en
bloques temporales contiguos; en cada partición un \textit{fold} distinto actúa como prueba y los
nueve restantes se usaron para entrenar y validar, como muestra la
figura~\ref{fig:within}.

\begin{figure}[H]
\centering
\begin{tikzpicture}[
  font=\footnotesize,
  bloque/.style={draw, minimum width=0.82cm, minimum height=0.75cm},
  fl/.style={-{Latex[length=1.8mm]}, semithick},
  x=1cm, y=1cm
]
\foreach \i in {1,...,10}{
  \node[bloque, fill=blue!8] at ({0.9*\i},0) {\scriptsize \i};
}
\node[bloque, fill=orange!25] (test) at ({0.9*3},0) {\scriptsize 3};

\node[anchor=south] (etq) at ({0.9*3},0.75) {\scriptsize prueba: 20 ensayos};
\draw[fl] (etq.south) -- (test.north);

\node[anchor=north, align=center] at (4.95,-0.55)
  {\scriptsize los otros nueve \textit{folds} (180 ensayos):
   144 de entrenamiento y 36 de validación, que determinan la parada temprana};
\node[anchor=north, align=center] at (4.95,-1.35)
  {\scriptsize el \textit{fold} de prueba rota $\rightarrow$ 10 particiones por sesión,
   2 sesiones por sujeto};
\end{tikzpicture}
\caption{Validación intra-sujeto. Los 200 ensayos de una sesión se dividen en
diez \textit{folds} estratificados por clase, sorteados al azar dentro de la sesión (no son
bloques temporales contiguos), y cada uno actúa como conjunto de prueba una vez. El proceso se
repite en las dos sesiones, lo que da 20 particiones por sujeto, y con tres semillas ($2026$, $2027$ y $2028$) en las
cuatro arquitecturas (sección~\ref{sec:semillas-por-arq}).}
\label{fig:within}
\end{figure}

Debe señalarse que en este protocolo cada modelo se entrenó con
\textbf{144 ensayos} y se detiene con 36. Es un régimen de datos muy escaso que, junto
con el reducido tamaño de la validación (sección~\ref{sec:seleccion-intra}), condiciona
buena parte de lo que se observa en el capítulo~\ref{cap:resultados}.

\subsection{Validación \loso{}}

Se emplea validación \textit{Leave-One-Subject-Out} (\loso{}) sobre los 54 sujetos. En cada
\textit{fold}, un sujeto constituye el conjunto de prueba y los $53$ restantes se reparten
en dos grupos: \textbf{dos sujetos completos} forman el conjunto de validación, cuya pérdida
determina la parada temprana, y los otros \textbf{51} forman el conjunto de entrenamiento. Los
dos sujetos de validación se eligen de forma determinista y rotan entre \textit{folds}, de
modo que ninguno actúa siempre como validación. El sujeto de prueba \textbf{no aporta ningún
ensayo} al entrenamiento, a la validación ni a los estadísticos de normalización.

Cada modelo se entrenó así con $51 \times 360 = 18\,360$ ensayos, frente a los 144 del
protocolo anterior, y se detiene con $2 \times 400 = 800$ ensayos de personas que no
participan en el ajuste de los pesos. Que la parada temprana se decida sobre personas no
vistas tiene una consecuencia que conviene anotar: el modelo conservado es el que mejor
generaliza a personas nuevas dentro del conjunto de origen, y no el que mejor se ajusta a
las personas con las que se entrenó.

\subsection{Protocolos adicionales: modelo global y modelo con sujetos apartados}
\label{sec:protocolos-extra}

Los dos protocolos anteriores responden a preguntas dentro de \openbmi{}. Para evaluar el
cambio de dominio entre conjuntos de datos distintos hacen falta dos protocolos más, que se
describen aquí porque condicionan la lectura de todos los resultados del
capítulo~\ref{cap:validacion-externa}.

\subsubsection{Qué sesiones se emplean}

Conviene fijarlo antes que nada, porque no es simétrico entre los dos conjuntos y de ello
depende la interpretación de las cifras:

\begin{itemize}
  \item De \openbmi{} se emplean \textbf{sus dos sesiones}, en todos los protocolos. Cada
  sujeto aporta $400$ segmentos, $200$ por sesión, y los $54$ sujetos suman $21\,600$.
  \item De WBCIC-SHU se emplea \textbf{únicamente la primera de sus tres sesiones}, tanto
  para entrenar el modelo nativo como para evaluar la transferencia. Cada sujeto aporta
  $200$ segmentos. La versión distribuida contiene $10\,199$, uno menos que los $10\,200$
  nominales, y quedan $10\,198$ tras descartar un segmento saturado
  (sección~\ref{sec:wbcic-preproc}).
\end{itemize}

Restringirse a la primera sesión de WBCIC-SHU es una elección conservadora: sus autores
documentan un efecto de aprendizaje entre días: con \eegnet{}, la exactitud media sube de
$0.818$ en la primera sesión a $0.866$ en la segunda y $0.889$ en la tercera
\cite{yang2025wbcic}, de modo que la primera es la
condición menos favorable. La consecuencia, que se declara aquí y se retoma al interpretar
los resultados, es que el modelo nativo de WBCIC entrena con $144$ segmentos por persona
frente a los $360$ de un modelo de \openbmi{}: una sesión menos y, además, la fracción que
cede al conjunto reservado.

\subsubsection{Protocolo global: un solo modelo sobre todos los sujetos}

En los protocolos intra-sujeto y \loso{} el resultado es un promedio sobre decenas de
modelos: $20$ particiones por sujeto y semilla en el primero y $54$ modelos distintos en el segundo,
cada uno entrenado con $51$ sujetos y detenido con otros $2$. Esto impide trasladar \emph{un} modelo a otro conjunto de
datos, porque no existe un único conjunto de pesos que represente al protocolo.

El protocolo global entrena, por tanto, \textbf{un solo modelo por arquitectura} con la
totalidad de los sujetos disponibles, reservando un $10$\,\% estratificado por sujeto, clase
y sesión para la parada temprana. No produce ninguna medida de generalización por sí mismo:
su única salida son los pesos que después se aplicaron al segundo conjunto de datos. La
figura~\ref{fig:global} lo esquematiza.

\begin{figure}[H]
\centering
\begin{tikzpicture}[
  font=\footnotesize,
  caja/.style={draw, rounded corners=2pt, minimum height=0.8cm, align=center},
  fl/.style={-{Latex[length=1.8mm]}, semithick},
  x=1cm, y=1cm
]
\node[caja, fill=blue!10, minimum width=7.2cm] (tr) at (0,0)
  {\textbf{Entrenamiento}: 54 sujetos, 2 sesiones\\45 canales, $19\,440$ segmentos ($90$\,\%)};
\node[caja, fill=blue!22, minimum width=2.4cm] (va) at (4.9,0)
  {\textbf{Validación}\\$2\,160$ ($10$\,\%)\\[-2pt]\scriptsize parada temprana};
\node[anchor=north, font=\scriptsize, align=center] at (0,-0.85)
  {los mismos 54 sujetos en las dos cajas: la validación mide\\
   \textbf{sujetos vistos}, no generalización};
\node[caja, fill=orange!20, minimum width=4.4cm] (ck) at (0,-2.7)
  {Un único conjunto de pesos por arquitectura};
\draw[fl] (tr.south) ++(0,-1.2) -- (ck.north);
\node[caja, fill=green!14, minimum width=5.6cm] (ap) at (0,-4.2)
  {Aplicación a WBCIC-SHU \textbf{sin reentrenar}};
\draw[fl] (ck) -- (ap);
\end{tikzpicture}
\caption{Protocolo global. Se entrenó un solo modelo por arquitectura con los 54 sujetos
para disponer de un conjunto de pesos transferible, con los 45 canales comunes a los dos
conjuntos. La validación determina la parada temprana, y su exactitud \textbf{no es una medida
de generalización}: procede de segmentos de los mismos sujetos que formaron parte del
entrenamiento.}
\label{fig:global}
\end{figure}

Este protocolo tiene una limitación que conviene enunciar de forma explícita, porque de
ella se deriva el siguiente: al entrenar con los $54$ sujetos \textbf{no queda ninguna persona
sin ver}. El modelo resultante no puede producir una cifra de \emph{persona nueva dentro
del mismo protocolo de adquisición}, y por tanto la caída que se observe al aplicarlo a otro
conjunto de datos mezcla dos cambios simultáneos, el de la persona y el del protocolo, sin
permitir separarlos.

Tampoco es lícito estimar el primero de esos dos efectos recurriendo al valor \loso{}. La
exactitud \loso{} es la media de $54$ modelos entrenados con $51$ sujetos cada uno, mientras
que la del protocolo global procede de un único modelo entrenado con $54$: son
\textbf{modelos distintos}, y su diferencia no aísla ningún factor.

\subsubsection{Protocolo con sujetos apartados: 44 más 10}

Para medir por separado el efecto de la persona y el del protocolo se entrenó una segunda
tanda de modelos apartando \textbf{diez sujetos completos} de \openbmi{}, que no aportan
ningún segmento ni al entrenamiento ni a los estadísticos de normalización. Un mismo
conjunto de pesos se evaluó entonces sobre tres conjuntos, como muestra la
figura~\ref{fig:hold}. Del primero al segundo cambia la persona. Del segundo al tercero
cambia un \textbf{factor compuesto}: a la vez la población, la institución, el amplificador,
el montaje y las diferencias de preprocesamiento de la tabla~\ref{tab:armonizacion}, que
este diseño no permite separar.

\begin{figure}[H]
\centering
\begin{tikzpicture}[
  font=\footnotesize,
  caja/.style={draw, rounded corners=2pt, minimum height=0.8cm, align=center},
  ev/.style={draw, rounded corners=2pt, fill=green!14, minimum width=3.5cm,
             minimum height=1.35cm, align=center, font=\scriptsize},
  fl/.style={-{Latex[length=1.8mm]}, semithick},
  x=1cm, y=1cm
]
\node[caja, fill=blue!10, minimum width=5.6cm] (tr) at (-1.3,0)
  {\textbf{44 sujetos} de \openbmi{}\\45 canales, $15\,840$ segmentos};
\node[caja, fill=blue!22, minimum width=2.0cm] (va) at (2.6,0)
  {\textbf{Val.}\\$1\,760$};
\node[caja, fill=red!12, minimum width=2.6cm] (ho) at (5.6,0)
  {\textbf{10 sujetos}\\apartados};
\node[anchor=north, font=\scriptsize] at (5.6,-0.62) {nunca vistos};

\node[caja, fill=orange!20, minimum width=4.0cm] (ck) at (0,-1.9)
  {Un conjunto de pesos};
\draw[fl] (tr.south) -- (ck.north);

\node[ev] (e1) at (-4.6,-4.0)
  {\textbf{1. Validación}\\sujetos vistos\\determina la parada temprana};
\node[ev] (e2) at (0,-4.0)
  {\textbf{2. Los 10 apartados}\\\textbf{persona nueva}\\protocolo conocido};
\node[ev] (e3) at (4.6,-4.0)
  {\textbf{3. WBCIC-SHU}\\persona nueva\\\textbf{protocolo nuevo}};
\coordinate (bus) at (0,-2.9);
\draw[fl] (ck) -- (bus);
\foreach \e in {e1,e2,e3} \draw[fl] (bus) -| (\e.north);

\draw[{Latex[length=1.6mm]}-{Latex[length=1.6mm]}, dashed] (e1.south) ++(0,-0.25)
  -- node[below, font=\scriptsize] {efecto de la persona} ($(e2.south)+(0,-0.25)$);
\draw[{Latex[length=1.6mm]}-{Latex[length=1.6mm]}, dashed] (e2.south) ++(0,-0.95)
  -- node[below, font=\scriptsize] {efecto del protocolo} ($(e3.south)+(0,-0.95)$);
\end{tikzpicture}
\caption{Protocolo con sujetos apartados. Los tres conjuntos de evaluación se recorren con
\textbf{los mismos pesos}, de modo que las diferencias entre ellos no se deben a modelos
distintos: de 1 a 2 cambia la persona; de 2 a 3 cambia además el protocolo de adquisición,
que agrupa varios factores (tabla~\ref{tab:armonizacion}).
El conjunto 1 es también el que determina la parada temprana, y sus ensayos proceden de las
mismas sesiones que los de entrenamiento; ambas condiciones tienden a elevar su exactitud
y, con ella, la caída de 1 a 2.}
\label{fig:hold}
\end{figure}

La selección de los diez sujetos es aleatoria con semilla fija y se recalcula de forma
determinista a partir de la lista ordenada de sujetos, sin depender del orden en que el
sistema de ficheros los devuelva. Los sujetos apartados resultan ser
sub-04, 07, 13, 17, 30, 34, 36, 47, 51 y 52. Su exactitud \loso{} media, con la media de tres semillas por sujeto, queda por debajo
de la del resto en las cuatro arquitecturas, entre $2.2$ y $3.1$ puntos: con \eegnet{},
$0.742$ frente a $0.765$. La diferencia no es significativa (Mann-Whitney, $p = 0.56$ con
\eegnet{} y entre $0.34$ y $0.58$ con las otras tres), pero con diez sujetos la prueba no
tiene potencia para descartar un déficit de ese tamaño, que coincide con el margen de
equivalencia de la sección~\ref{sec:equivalencia}. Como el efecto de la persona se calcula
restando la exactitud sobre los apartados a la de los sujetos vistos, esa selección puede
sobrestimarlo en una magnitud comparable. Se hace
constar que \textbf{la semilla no se volvió a sortear tras comprobar esto}: repetir el
sorteo hasta obtener un subconjunto favorable sería una selección sesgada.

El diseño original empleó un único sorteo de diez sujetos, por lo que la estimación del \idp{} fue imprecisa. Se planteó una extensión a cinco particiones deterministas de $11$, $11$, $11$, $11$ y $10$ sujetos, con semilla $2026$ y asignaciones registradas en el archivo de particiones del repositorio, para que cada uno de los $54$ sujetos actuara una vez como persona nueva. Cada partición se entrenó con los sujetos restantes y conservó el esquema y la validación interna originales, con ensayos reservados de los sujetos de entrenamiento. Se previó ejecutar \eegnet{} con tres semillas en local y las cuatro arquitecturas en el laboratorio, mediante los scripts de ejecución y análisis de las cinco particiones.
% PENDIENTE-E2/E3: resultados

\subsubsection{Normalización del conjunto de destino}

Al aplicar un modelo a un conjunto de datos que no formó parte de su entrenamiento, la media y la desviación con que se
normaliza la entrada no están definidas de forma única. Se reportaron tres alternativas, porque la
diferencia entre ellas es en sí misma informativa sobre la naturaleza del cambio de dominio:

\begin{itemize}
  \item \textbf{Origen}: se aplicaron la media y la desviación del entrenamiento en
  \openbmi{}. Es la condición más estricta, en la que no se emplea ninguna información del
  conjunto de destino.
  \item \textbf{Por segmento}: cada segmento se escala con una única media y una única
  desviación, calculadas sobre todos sus canales. No emplea ningún otro segmento, de modo
  que no es transductiva, y corrige una diferencia de ganancia común a todos los canales;
  no corrige diferencias de ganancia entre canales ni de covarianza espacial.
  \item \textbf{Por sujeto}: estadísticos de cada sujeto del destino. Sí es transductiva,
  pues supone disponer de todos los datos del sujeto antes de clasificar.
\end{itemize}

Esta elección permite contrastar una hipótesis concreta y acotada: \textbf{si la caída al
cambiar de conjunto se debiera a una diferencia de ganancia global entre los equipos de
registro, la normalización por segmento la corregiría y el desempeño aumentaría}. El
contraste tiene dos límites. No alcanza a diferencias de ganancia por canal ni de
covarianza. Y las tres políticas se aplicaron a un modelo entrenado con los estadísticos por
canal del origen, de modo que las dos alternativas introducen además un desajuste entre
la normalización de entrenamiento y la de prueba; un contraste completo exigiría entrenar
también con cada normalización. El capítulo~\ref{cap:validacion-externa} muestra que la
normalización por segmento no recupera el desempeño.

Finalmente, para que el modelo transferido y un modelo entrenado en el destino sean
comparables, se aparta en WBCIC-SHU un $20$\,\% de los segmentos de \textbf{cada} sujeto,
estratificado por sujeto y clase, con una semilla propia e independiente de la del
entrenamiento. Ese subconjunto, $2\,040$ segmentos, exactamente $40$ de cada uno de los
$51$ sujetos y equilibrado entre clases, no se emplea para entrenar ni para la parada
temprana de ninguno de los dos modelos, y es donde se miden ambos.

\subsection{Calibración con pocos ensayos del sujeto no visto}
\label{sec:metodo-fewshot}

El protocolo \loso{} mide el desempeño con cero ensayos del sujeto de prueba, y el
intra-sujeto con $144$ de entrenamiento. Para el escenario intermedio se calibra el modelo
con $k$ ensayos del sujeto excluido, con $k \in \{10, 20\}$. La calibración se ejecuta
dentro de cada \textit{fold} \loso{}, justo después de evaluar el modelo ya entrenado, sin
reentrenarlo, de modo que cada \textit{fold} y semilla tiene un único modelo de partida.

Los $400$ ensayos del sujeto excluido, de sus dos sesiones juntas, se dividen así. Para cada
$k$ y cada uno de $5$ sorteos se eligen $k$ ensayos de calibración, estratificados por clase
($k/2$ por clase), con un generador inicializado con $\mathrm{semilla}\times 1000 + k\times 10 +
\mathrm{sorteo}$. La evaluación es el resto del sujeto ($390$ ensayos si $k = 10$, $380$ si
$k = 20$); ningún ensayo de calibración se evaluó. Se compararon cuatro métodos, uno de ellos (cero) la referencia, sobre el mismo
sorteo y el mismo resto, de modo que son pareados:

\begin{itemize}
  \item \textbf{Cero.} El modelo \loso{} sin calibrar, medido sobre ese resto. Es la
  referencia pareada: sin ella, la comparación mezclaría el efecto de la calibración con el
  de evaluar sobre menos ensayos.
  \item \textbf{Estadísticos} (\texttt{stats}). Se sustituyen la media y la desviación por canal del
  entrenamiento por las de los $k$ ensayos de calibración. No se modifica ningún peso.
  \item \textbf{Cabeza} (\texttt{cabeza}). Se reentrena solo la última capa de clasificación de la
  arquitectura con los $k$ ensayos, y el resto del modelo permanece congelado. Se emplea
  AdamW con tasa de aprendizaje $10^{-3}$ y decaimiento de pesos $10^{-2}$, entropía cruzada
  sin suavizado y $100$ pasos de lote completo, con las capas de normalización por lotes y
  el \textit{dropout} en modo de evaluación (recalcular sus estadísticos con $10$ ensayos
  sería otra forma de adaptación, no la que se quiere medir). La semilla de la cabeza es la
  semilla de la ejecución más el índice del sorteo.
  \item \textbf{Estadísticos y cabeza} (\texttt{stats+cabeza}). Las dos operaciones: la cabeza se entrenó sobre los
  ensayos de calibración ya normalizados con sus propios estadísticos.
\end{itemize}

El número de pasos y la tasa de aprendizaje de la cabeza se fijaron de antemano y no se
ajustaron en ningún conjunto; son una elección del autor, sin una referencia que los respalde, y el
resultado vale para esa configuración. No hay conjunto de validación ni parada temprana en la
calibración. La \emph{ganancia} de un método es su exactitud menos la de \emph{cero} sobre
el mismo sorteo; se promedia entre los cinco sorteos dentro de cada sujeto, y entre
semillas dentro de él.

La calibración se ejecutó sobre los modelos \loso{} de las semillas $2027$ y $2028$ y sobre
los de la variante con alineamiento (semilla $2026$).


\subsection{Indicador de degradación por cambio de dominio}
\label{sec:dsi}

\subsubsection{Problema del indicador clásico}

El indicador de uso común cuantifica la degradación al pasar de la validación
intra-sujeto (\textit{subject-dependent}) a la condición inter-sujeto
(\textit{subject-independent}), comparación habitual en la literatura de interfaces cerebro-computadora
(BCI) independiente del sujeto \cite{kwon2020}. Esa literatura compara ambos protocolos; la
forma relativa que sigue es la formalización empleada en este trabajo para contrastarla:

\begin{equation}
\mathrm{DSI}_{\text{clásico}} =
  \frac{m_{\text{intra}} - m_{\text{\loso{}}}}{m_{\text{intra}}}
\label{eq:dsi-clasico}
\end{equation}

donde $m$ es una métrica de desempeño. Esta formulación es \textbf{confundida}:
las dos condiciones no difieren solo en si el sujeto de prueba fue visto, sino
también en el tamaño del conjunto de entrenamiento. En \openbmi{}, la condición
intra-sujeto entrena con unos $10^2$ ensayos y la \loso{} con $10^4$. La
diferencia de desempeño resultante mezcla el efecto del cambio de dominio con el
efecto del volumen de datos, y no permite aislar el primero. La
sección~\ref{sec:within-vs-loso} cuantifica esa confusión sobre los datos.

\subsubsection{Indicador propuesto: \texorpdfstring{\dsi{}}{IDS}, índice de degradación por sujeto}

Para evitar la confusión con el indicador clásico, que conserva la sigla \textit{DSI}
establecida en la literatura, el propuesto aquí se denomina \textbf{\dsi{}, índice de
degradación por sujeto}. Se propone comparar, \emph{para un mismo modelo entrenado una sola
vez}, el desempeño sobre dos conjuntos de prueba que difieren únicamente en si el sujeto
que los produjo estuvo o no en el entrenamiento, como se definió en la ecuación~\eqref{eq:dsi}:

\begin{equation}
\mathrm{IDS} = \frac{m_{\text{visto}} - m_{\text{no visto}}}{m_{\text{visto}}}
\label{eq:dsi}
\end{equation}

donde:

\begin{itemize}
  \item $m_{\text{no visto}}$ se evaluó sobre el sujeto excluido de la partición
        \loso{};
  \item $m_{\text{visto}}$ se evaluó sobre una fracción de ensayos
        \textbf{reservada} de los 51 sujetos de entrenamiento, que no participa
        en el ajuste de los pesos, en el cálculo de los estadísticos de
        normalización ni en la parada temprana. Esta última la determina un conjunto
        distinto, el de los dos sujetos de validación, de modo que $m_{\text{visto}}$
        no queda favorecido por el criterio de selección del modelo.
\end{itemize}

La reserva se estratifica por sujeto, clase y sesión, de modo que no se
introduzca cambio temporal (día 1 frente a día 2) en la comparación. Se fijó una
fracción del 10\,\%.

Al compartir modelo, pesos y datos de entrenamiento, ambos conjuntos difieren en la
pertenencia del participante al entrenamiento, y el indicador queda así libre del
confundidor de volumen de datos que afecta a la ecuación~\eqref{eq:dsi-clasico}. No es, sin
embargo, la única diferencia entre ellos. Los ensayos vistos se reservan de las mismas
sesiones que los de entrenamiento, de modo que el modelo se ha ajustado con datos no solo
de esa persona sino también de esa misma sesión de registro, y el indicador refleja ambas
condiciones. Además, el conjunto visto promedia 51 personas y el no visto corresponde a una sola.
Por ello, el indicador de cada \textit{fold} depende casi por completo de la dificultad
del participante excluido, y solo se interpreta de forma agregada. La figura~\ref{fig:particiones} resume el
esquema.

\begin{figure}[H]
\centering
\begin{tikzpicture}[
  font=\footnotesize,
  caja/.style={draw, rounded corners=2pt, minimum height=0.85cm, align=center},
  fl/.style={-{Latex[length=2mm]}, semithick},
  ev/.style={-{Latex[length=2mm]}, semithick, densely dashed, gray!70!black},
  x=1cm, y=1cm
]
\node[caja, fill=black!5, minimum width=9.4cm] (todos) at (0,0)
  {54 sujetos del conjunto de datos};

\node[caja, fill=blue!8,   minimum width=4.6cm] (tr) at (-2.4,-1.7)
  {51 sujetos de entrenamiento};
\node[caja, fill=blue!20,  minimum width=2.0cm] (va) at ( 2.0,-1.7)
  {2 sujetos\\[-3pt] \scriptsize parada temprana};
\node[caja, fill=orange!25, minimum width=2.2cm] (te) at ( 4.6,-1.7)
  {sujeto $s$\\[-3pt] \scriptsize \textbf{no visto}};
\draw[fl] (todos.south -| tr.north) -- (tr.north);
\draw[fl] (todos.south -| va.north) -- (va.north);
\draw[fl] (todos.south -| te.north) -- (te.north);

\node[caja, fill=blue!8,  minimum width=2.6cm] (entren) at (-3.4,-3.4)
  {entrenamiento\\[-3pt] \scriptsize 90\,\% de los ensayos};
\node[caja, fill=green!20, minimum width=1.8cm] (visto) at (-1.0,-3.4)
  {\textbf{visto}\\[-3pt] \scriptsize 10\,\%};
\draw[fl] (tr.south -| entren.north) -- (entren.north);
\draw[fl] (tr.south -| visto.north) -- (visto.north);

\node[caja, fill=black!10, minimum width=2.6cm] (modelo) at (-3.4,-5.1)
  {modelo\\[-3pt] \scriptsize entrenado una vez};
\draw[fl] (entren) -- (modelo);

\draw[ev] (modelo.south) -- (-3.4,-6.2) -- (-1.0,-6.2)
      node[midway, below, font=\scriptsize] {evalúa} -- (visto.south);
\draw[ev] (modelo.south) -- (-3.4,-6.9) -- (4.6,-6.9)
      node[midway, below, font=\scriptsize] {evalúa} -- (te.south);

\node[font=\scriptsize, anchor=west] at (-0.05,-3.4) {\;$\rightarrow m_{\text{visto}}$};
\node[font=\scriptsize, anchor=west] at ( 5.7,-1.7) {\;$\rightarrow m_{\text{no visto}}$};

\node[caja, draw=black, thick, fill=yellow!15, minimum width=6.4cm] (dsi) at (0,-8.3)
  {$\mathrm{IDS} = \dfrac{m_{\text{visto}} - m_{\text{no visto}}}{m_{\text{visto}}}$};
\end{tikzpicture}
\caption{Esquema de particiones para el cálculo del indicador de degradación.
En cada uno de los 54 \textit{folds} \loso{}, un sujeto queda excluido por completo
(\emph{no visto}), dos sujetos completos determinan la parada temprana y se reserva el
10\,\% de los ensayos de los 51 de entrenamiento (\emph{visto}), estratificado por
sujeto, clase y sesión. Ambos conjuntos se evaluaron con el mismo modelo y los mismos
pesos, y ninguno de los dos interviene en la selección del modelo.}
\label{fig:particiones}
\end{figure}

\subsection{Criterio operativo de robustez}
\label{sec:criterio-robustez}

El objetivo general de este trabajo plantea la identificación de la arquitectura de mayor
robustez ante el cambio de dominio, lo que exige una definición operativa del término. La
exactitud promedio no la proporciona: las cuatro arquitecturas quedan comprendidas en un
intervalo de $4.6$ puntos en el protocolo \loso{}, y los contrastes pareados de la
sección~\ref{sec:estadistica} no sostienen una diferencia entre las tres primeras.
Calificar como más robusta a la de media más alta equivaldría a atribuir significado a una
diferencia que los datos no respaldan.

Se define aquí la robustez como la estabilidad del desempeño ante desplazamientos del
dominio, y se evaluó sobre tres ejes independientes, medidos con los protocolos descritos
en las secciones precedentes.

El primero es la \textbf{degradación ante participantes no vistos}, cuantificada mediante
el \dsi{} de la sección~\ref{sec:dsi}, que compara un mismo conjunto de pesos sobre
participantes vistos y no vistos. Valores menores corresponden a mayor robustez, con una
salvedad que afecta también al segundo eje: un clasificador en el nivel del azar tendría un
indicador nulo sin ser robusto, de modo que ambos índices se interpretan junto con la
exactitud sobre participantes no vistos.

El segundo es la \textbf{degradación ante un protocolo de adquisición distinto}. Se
introduce para ello un indicador análogo al anterior, el \textbf{\idp{}, índice de
degradación por protocolo}, construido sobre la misma lógica de un único conjunto de pesos
evaluado en dos conjuntos de prueba, según la ecuación~\eqref{eq:idp}:

\begin{equation}
\mathrm{IDP} =
  \frac{m_{\text{participante nuevo}} - m_{\text{protocolo nuevo}}}
       {m_{\text{participante nuevo}}}
\label{eq:idp}
\end{equation}

donde $m_{\text{participante nuevo}}$ corresponde a participantes no vistos del conjunto de
origen y $m_{\text{protocolo nuevo}}$ a participantes del conjunto externo. Ambos términos
comparten la condición de no haber sido vistos durante el entrenamiento, de modo que la
diferencia entre ellos no vuelve a contabilizar el cambio de persona que el \dsi{} ya recoge.
Lo que mide es el cambio de protocolo de adquisición en sentido amplio, un factor compuesto
de los elementos de la tabla~\ref{tab:armonizacion}, y no un factor aislado. Su cálculo exige el protocolo de sujetos apartados descrito
en la sección~\ref{sec:protocolos-extra}, pues un modelo entrenado sobre la totalidad de los
participantes no deja ninguno sin ver y no puede producir el primero de los dos términos.
Los dos índices son adimensionales, pero no proceden del mismo régimen: el \dsi{} se calcula
en \loso{} con $51$ sujetos de entrenamiento y $62$ canales, y el \idp{} con un único modelo
de $44$ sujetos y $45$ canales. La comparación de sus magnitudes es, por tanto, aproximada.

El tercero es la \textbf{estabilidad entre participantes}. La exactitud promedio no informa
sobre cómo se reparte el desempeño entre las personas, reparto que resulta determinante
para una interfaz destinada a usuarios arbitrarios. Se cuantifica mediante la desviación
típica de la exactitud por participante, el percentil $5$ de esa distribución y el número
de participantes situados por debajo del umbral funcional del $70$\,\%.

Los tres ejes se reportan de forma conjunta y no se agregan en un índice único. La
agregación requeriría fijar pesos relativos que ninguna consideración del problema
justifica, y encubriría el hecho, relevante en sí mismo, de que las arquitecturas se
ordenan de manera distinta según el eje considerado.

La interpretación del tercer eje admite una ambigüedad que conviene resolver de antemano.
Una arquitectura cuya exactitud media se aproxima al azar comprime necesariamente el rango
de su distribución por participante, de modo que una desviación baja resulta compatible
tanto con un desempeño estable como con una capacidad discriminativa reducida y uniforme.
El percentil $5$ y el recuento bajo umbral separan ambas situaciones, razón por la cual el
eje no se sustenta en la desviación de forma exclusiva. A ello se añade que los tres ejes
se estiman con tres semillas por arquitectura (sección~\ref{sec:semillas-por-arq}), cuya
desviación, calculada con tres valores, es una estimación gruesa, y que el régimen de
optimización no es homogéneo entre esas semillas en \eegnet{}, el Conformer y LaBraM
(sección~\ref{sec:recorte-gradiente}); por ello las diferencias de pequeña magnitud entre
arquitecturas conservan carácter indicativo.

\subsection{Análisis estadístico}
\label{sec:estadistica}

\subsubsection{Métricas}

Se reportan dos. La \textbf{exactitud} es la proporción de ensayos clasificados
correctamente; con dos clases equiprobables, el azar está en 0.50. El
\textbf{coeficiente kappa} \cite{cohen1960} mide lo mismo pero descontando los aciertos
esperables por azar: vale 0 cuando el modelo no aporta nada sobre adivinar y 1
cuando acierta todo. Con dos clases balanceadas en el conjunto de prueba, como es el caso,
kappa es una transformación lineal exacta de la exactitud, $\kappa = 2\,\mathrm{exactitud}-1$,
de modo que ambas métricas ordenan igual y producen los mismos contrastes; que coincidan no
constituye una corroboración. Se reporta kappa porque el \dsi{} calculado sobre él mide la
degradación relativa al exceso sobre el azar, que no depende del nivel base de $0.5$.

\subsubsection{Comparación entre arquitecturas}

Todas las arquitecturas se evaluaron sobre \textbf{las mismas 54 particiones}, con
los mismos sujetos en el mismo orden. Esto permite compararlas sujeto a sujeto en
lugar de comparar solo los promedios, que es más informativo: la variabilidad
entre personas es grande y podría impedir detectar una diferencia real entre modelos.

Para decidir si una diferencia observada es atribuible al azar se emplea la
\textbf{prueba de Wilcoxon para muestras pareadas} \cite{wilcoxon1945}, que compara los
resultados sujeto por sujeto sin suponer que sigan una distribución normal. Es bilateral,
sobre los $54$ pares, y reparte las diferencias nulas entre ambos signos
(\texttt{zero\_method="zsplit"} de SciPy). Se acompaña
siempre del número de sujetos en que cada arquitectura resulta mejor, que es una
lectura directa de la consistencia del efecto, y del tamaño de efecto
$r = |Z|/\sqrt{n}$ \cite{fritz2012}, con $Z = \Phi^{-1}(1 - p/2)$ reconstruido a partir del
valor $p$ bilateral y $n = 54$. Las cuatro arquitecturas cuentan con tres
semillas en los dos protocolos. Los contrastes se calcularon sobre la media de las tres semillas
($2026$, $2027$ y $2028$) de cada participante, de modo que cada participante aporta un valor por
arquitectura, el número de pares es siempre $54$ y las semillas no se tratan como réplicas
independientes (sección~\ref{sec:estadistica-variantes}); los resultados con la
semilla $2026$ sola y por semilla se informan en la
sección~\ref{sec:variabilidad-semillas}.

Cada protocolo y cada métrica constituyen una familia de seis contrastes, uno por par de
arquitecturas, fijada antes de calcular las cifras con tres semillas. Se reportaron los valores $p$
sin corregir y corregidos por Holm \cite{holm1979} dentro de cada familia; con tres semillas la
corrección no cambia ninguna conclusión de significación. La corrección de Holm ordena
los $m$ valores $p$ de menor a mayor, multiplica el $i$-ésimo por $m - i + 1$ e impone que
la secuencia ajustada no decrezca.

\subsubsection{Equivalencia práctica}
\label{sec:equivalencia}

Que una diferencia no alcance significación no demuestra que sea pequeña, ni tampoco lo
demuestra un tamaño de efecto observado reducido. Para sostener que dos arquitecturas rinden
igual se emplea un criterio de equivalencia \cite{schuirmann1987,lakens2017}: se fija un
margen $\delta$, se plantea como hipótesis nula $|\Delta| \geq \delta$ frente a la
alternativa $|\Delta| < \delta$, y se declaran
equivalentes cuando el intervalo de confianza al $90$\,\% de su diferencia pareada queda
íntegramente dentro de $(-\delta,\ +\delta)$, lo que corresponde a dos contrastes unilaterales
al $5$\,\%. El intervalo se obtuvo por remuestreo pareado de percentiles
(sección~\ref{sec:bootstrap}) y se aplicó a la exactitud por participante, en \loso{} y en
intra-sujeto, calculada como media de las tres semillas de cada participante, de modo que la
variabilidad por inicialización se promedia dentro del participante y no queda incorporada en el
intervalo. Como comprobación se informan también el intervalo jerárquico de sujetos y semillas
(sección~\ref{sec:variabilidad-metodo}), que sí la incorpora, y el contraste de dos pruebas unilaterales (TOST, \textit{two one-sided tests}) con la prueba de
Wilcoxon (corrección de Holm en la familia de seis) y con la $t$ de Student.

Se fija $\delta = 0.03$ de exactitud. Equivale a entre un quinto y un tercio de la
desviación entre participantes, de $0.10$ a $0.15$ según arquitectura y protocolo, y es del
orden de la variabilidad entre semillas de la exactitud media de un participante, de
$0.016$ a $0.025$ en el protocolo intra-sujeto y de $0.019$ a $0.025$ en \loso{}, con las cuatro
arquitecturas y tres semillas. El margen se
fijó después de un primer análisis de los datos y no de forma previa. Por ese motivo, la
sección~\ref{sec:wilcoxon} informa también el menor margen con el que se cumpliría la
equivalencia.

\textbf{Criterios de equivalencia informados.} Con el margen $\delta = 0.03$ de la
sección~\ref{sec:equivalencia} se informan tres criterios, para que la conclusión no
dependa de uno solo: el intervalo al $90$\,\% por remuestreo pareado dentro de
$(-\delta, +\delta)$; la prueba de dos contrastes unilaterales con Wilcoxon, corregida por
Holm en la familia de seis; y la misma con la prueba $t$, sin corregir. Se aplicó también al
protocolo intra-sujeto. No se interpreta un resultado no significativo como ausencia de
efecto: se informa como que no se detecta diferencia y, si procede, con el resultado de
equivalencia; y un resultado significativo se acompaña de la magnitud.


\subsubsection{Prueba bayesiana de rangos con signo}

Se aplicó la prueba bayesiana de rangos con signo de \cite{benavoli2017} mediante \texttt{baycomp} 1.0.3, con $50\,000$ muestras y semilla $2026$. La entrada fueron las exactitudes pareadas por sujeto, promediadas entre las tres semillas, para cada par de arquitecturas en el mismo protocolo; también se comparó la degradación por sujeto, para la cual una caída menor es mejor. Este promedio incluye los cambios de régimen entre semillas descritos en la sección~\ref{sec:semillas-metodo}, por lo que no constituye una comparación restringida a un único régimen. Se fijó una región de equivalencia práctica (region of practical equivalence, ROPE) de $\pm 0.03$, igual al margen del TOST, con $0.01$ y $0.02$ como análisis de sensibilidad. La salida comprende $P(A\text{ mejor})$, $P(\text{equivalentes})$ y $P(B\text{ mejor})$; se declara la decisión correspondiente si una probabilidad alcanza $0.95$ y, de lo contrario, se informa como indecisa.

Frente a Wilcoxon, esta prueba distribuye la probabilidad posterior entre ventajas y equivalencia práctica, en lugar de limitarse a un valor $p$ de contraste. Frente al TOST, proporciona directamente una probabilidad de equivalencia dentro de la ROPE, en lugar de una decisión basada en valores $p$.

\subsubsection{Intervalos de confianza por remuestreo}
\label{sec:bootstrap}

El contraste de Wilcoxon indica si una diferencia es atribuible al azar, pero no cuánta
incertidumbre acompaña a cada cifra reportada. Para acotarla se emplea el
\textbf{remuestreo con reemplazo} (\textit{bootstrap}) \cite{efron1979} sobre los valores por participante
que el entrenamiento ya produjo, procedimiento que no requiere entrenar ningún modelo
adicional.

De los $54$ valores por participante de una arquitectura se extraen, con reemplazo, $54$
valores; se calculó la media de esa extracción; y el procedimiento se repitió
$B = 10\,000$ veces. El intervalo de confianza al $95$\,\% son los percentiles $2.5$ y
$97.5$ de esa distribución, sin corrección de sesgo. La semilla del remuestreo se fija en
$2026$, de modo que dos ejecuciones producen cifras idénticas; como cada cantidad consume
el generador en un orden distinto, la referencia para reproducir una cifra concreta es la
función del verificador que la calcula. El remuestreo trata los \textit{folds} \loso{} como
independientes, aunque comparten $49$ de sus $51$ sujetos de entrenamiento.

La aplicación del procedimiento al \dsi{} requiere una precisión adicional. El remuestreo
se realiza sobre los \textit{folds} y el indicador se recalcula en cada extracción con la
ecuación~\eqref{eq:dsi} aplicada a las medias, que es la forma en que se reporta; promediar
el indicador de cada \textit{fold} estimaría una cantidad distinta. En las diferencias
entre arquitecturas, por su parte, la extracción de participantes es común a las dos
arquitecturas del par, pues remuestrear cada una por separado destruiría la correlación
entre ellas e inflaría el intervalo de forma artificial.

Con tres semillas se añade un remuestreo jerárquico. En cada una de las $B = 10\,000$ repeticiones
se extraen con reemplazo $54$ participantes, dentro de cada uno se extraen con reemplazo sus tres
semillas, con índices independientes entre arquitecturas, y se promedian las semillas antes de
calcular la diferencia o el indicador. Remuestrear en los dos niveles cuenta dos veces el ruido de
semilla, que ya forma parte de la dispersión entre participantes, de modo que estos intervalos son
algo conservadores; con solo tres semillas son además una aproximación tosca de la variación por
inicialización. Los intervalos de remuestreo solo de participantes difieren de los jerárquicos en
menos de $0.003$ en las diferencias entre arquitecturas.

Los intervalos no incorporan corrección por comparaciones múltiples. Se reportan como
descripción de la incertidumbre de cada cifra; las afirmaciones sobre diferencias entre
arquitecturas se sustentan en el contraste pareado, con la comprobación de Holm descrita
arriba.

\subsubsection{Criterio para declarar degradación}

Se declara degradación por cambio de dominio únicamente si se cumplen dos
condiciones: que el \dsi{} sea positivo y que la diferencia entre el conjunto
visto y el no visto no sea atribuible al azar, según un contraste de Wilcoxon bilateral
entre la exactitud en vistos y en no vistos sobre los $54$ \textit{folds}, con
$\alpha = 0.05$. Si solo se cumple la primera, se
reporta el valor indicando expresamente que no es concluyente.

Estas dos condiciones se fijaron antes de ejecutar las arquitecturas restantes. La
corrección de Holm sobre los cuatro contrastes, uno por arquitectura, se incorporó
\textbf{después} de conocer los resultados de las cuatro, y determina la clasificación de una
de ellas: sin corrección, EEGMamba ($p = 0.038$) también se declararía degradada. Por eso
la sección~\ref{sec:dsi-resultados} informa ambos valores.

\subsubsection{Contrastes de las variantes y familias de comparaciones múltiples}
\label{sec:estadistica-variantes}

\textbf{Unidad y contraste.} La unidad de análisis es el sujeto ($n = 54$). En el protocolo
intra-sujeto, la exactitud de un sujeto en una semilla es la media de sus $20$ particiones;
en \loso{} es la de su \textit{fold}. Cuando hay varias semillas, se promedian dentro del sujeto o
se modelan como factor anidado; nunca se tratan como réplicas independientes. El contraste
de una variante es la diferencia por sujeto entre la variante y su referencia de la misma
semilla. Para cada contraste se informa la media y la mediana de la diferencia, el
intervalo al $95$\,\% por remuestreo de sujetos (percentiles, $B = 10\,000$, con un
generador \texttt{numpy.random.default\_rng(2026)} nuevo para cada contraste), la prueba de
Wilcoxon bilateral con reparto de las diferencias nulas, el tamaño de efecto $r$ y el número
de sujetos que mejoran y empeoran, como en los contrastes entre arquitecturas descritos arriba. Dado que kappa
es una transformación lineal de la exactitud, se contrasta la exactitud.

\textbf{Familias de comparaciones múltiples.} La corrección de Holm se aplica dentro de
cada familia, definida antes de calcular los contrastes de las variantes y de las semillas adicionales (resumen en la tabla~\ref{tab:familias}):

\begin{itemize}
  \item entre arquitecturas, seis contrastes por protocolo, métrica y variante de semilla
  (semilla $2026$, media de tres semillas, media de $2027$ y $2028$);
  \item ablaciones intra-sujeto, una familia por arquitectura: ocho contrastes para
  \eegnet{} y para el Conformer, y tres para LaBraM y para EEGMamba;
  \item \loso{} con alineamiento: cuatro contrastes, uno por arquitectura; \loso{} con
  reflexión de canales: cuatro, uno por arquitectura;
  \item calibración con pocos ensayos: veinticuatro contrastes de ganancia frente a cero
  (cuatro arquitecturas, dos valores de $k$, tres métodos), para los modelos base y,
  aparte, otros veinticuatro para los de la variante con alineamiento; el contraste de
  equivalencia de cada una forma una familia propia de veinticuatro;
  \item comparación de las políticas de normalización del destino: ocho contrastes;
  \item alineamiento en la transferencia: doce contrastes (tres evaluaciones por cuatro arquitecturas).
\end{itemize}

Como lectura secundaria declarada se aplicó Holm sobre los $22$ contrastes de las
ablaciones a la vez y sobre los $8$ de \loso{} con alineamiento y con reflexión a la vez. Como
análisis de sensibilidad, se repiten las familias de alineamiento, reflexión y las
ablaciones de LaBraM y EEGMamba con la referencia sustituida por la media por sujeto de las
semillas $2027$ y $2028$ de la base, entrenadas con recorte; para \eegnet{} y el Conformer en
intra-sujeto no existe esa referencia, porque sus semillas $2027$ y $2028$ se entrenaron sin
recorte (tabla~\ref{tab:regimen}). Todo análisis fuera de estas familias es exploratorio y no se
corrige.


\subsubsection{Variabilidad por semilla de las diferencias}
\label{sec:variabilidad-metodo}

\textbf{Variabilidad por semilla de las diferencias.} Para cada contraste entre
arquitecturas se calcula la diferencia media por semilla (emparejando semillas por índice) y
la desviación típica de esos tres valores, con la salvedad de que tres valores dan dos
grados de libertad. Además, el intervalo de la diferencia de las medias de tres semillas
se obtiene con el remuestreo jerárquico descrito en la sección~\ref{sec:bootstrap}, con un
generador \texttt{default\_rng(2026)} nuevo por contraste. El mismo procedimiento, con los \textit{folds} \loso{}
como unidad en el \dsi{} y con los sujetos apartados y los de WBCIC-SHU como unidades del
\idp{}, da los intervalos jerárquicos de ambos índices.


\subsection{Criterios adicionales al desempeño}
\label{sec:criterios-adicionales}

\textbf{Costo computacional.} Se registraron parámetros, operaciones de punto flotante (FLOP) por ensayo con \texttt{torch.utils.flop\_counter}, latencia mediana y percentil 95 en unidad central de procesamiento (CPU) y CUDA con lotes de $1$ y $64$, memoria pico CUDA con lote $64$ y tamaño del punto de control. El contador cubre solo operaciones compatibles. El costo de entrenamiento se calculó como suma de tiempos registrados por ejecución, no como tiempo de pared (\texttt{scripts/informe\_costo.py}); las mediciones de inferencia de EEGMamba no estuvieron disponibles por su dependencia de \texttt{mamba\_ssm}.

\textbf{Estabilidad entre semillas.} A partir de las exactitudes por sujeto en intra-sujeto y \loso{}, se calculó la desviación estándar entre medias por semilla, la media de tau de Kendall entre rankings de sujetos y el porcentaje de sujetos con signo consistente de la diferencia frente a \eegnet{} (\texttt{scripts/informe\_estabilidad.py}). Se examinaron las tres semillas y, por separado, $2027$ y $2028$, teniendo en cuenta los cambios de régimen documentados.

\textbf{Robustez a perturbaciones.} Con los pesos del protocolo $44$ más $10$ se evaluaron los diez sujetos apartados ante ruido gaussiano según relación señal-ruido (SNR), interferencia de red de $60$ Hz, anulación de canales, escala de amplitud y desplazamiento de ventana. Se emplearon cinco sorteos para condiciones aleatorias; escala y desplazamiento son deterministas. Las perturbaciones precedieron la normalización con estadísticas de origen y se calcularon intervalos de confianza por remuestreo de sujetos (\texttt{scripts/informe\_robustez.py}).

\textbf{Interpretabilidad.} Se representaron los filtros espaciales de \eegnet{} en topomapas y se calculó la saliencia como gradiente absoluto del logit de la clase real respecto a la entrada normalizada, promediado en tiempo y ensayos por clase. Para la saliencia $S$ se definieron los índices de lateralización de la ecuación~\eqref{eq:lateralizacion}:
\begin{equation}
\label{eq:lateralizacion}
L_{C3,C4}=\frac{S_{C4}-S_{C3}}{S_{C4}+S_{C3}},\qquad
L_{\mathrm{vec}}=\frac{\overline S_{\mathrm{der}}-\overline S_{\mathrm{izq}}}{\overline S_{\mathrm{der}}+\overline S_{\mathrm{izq}}},
\end{equation}
donde las medias izquierda y derecha comprenden FC3/C3/CP3 y FC4/C4/CP4, respectivamente (\texttt{scripts/informe\_interpretabilidad.py}).

Robustez e interpretabilidad se limitaron a los diez apartados y una semilla de entrenamiento ($2026$); EEGMamba quedó fuera de ambos análisis porque \texttt{mamba\_ssm} no estaba disponible en el entorno de inferencia Windows. La eficiencia de datos se examinó mediante la calibración con pocos ensayos de la sección~\ref{sec:metodo-fewshot}.

\textbf{Regla de decisión.} Entre arquitecturas equivalentes tanto en exactitud \loso{} como en degradación, según TOST con margen $\pm 0.03$ y la prueba bayesiana con $P(\text{equivalentes})\geq 0.95$, se prefirió la de menor costo computacional. Esta regla se fijó después del primer análisis de desempeño, a pedido de la dirección, y no antes de las ejecuciones.

\subsection{Control de fuga de datos}

El particionado \loso{}, el intra-sujeto y el conjunto de vistos están cubiertos por
\textbf{41 pruebas automatizadas}: 22
sobre el particionado \loso{} y la separación de sujetos apartados, 12 sobre el
intra-sujeto y 7 sobre el conjunto reservado. No las tienen, en cambio, las funciones que
construyen el protocolo global, el de sujetos apartados y el $20$\,\% reservado de
WBCIC-SHU; para esos tres se comprobó que los recuentos de segmentos coinciden con los
esperados. Las pruebas verifican, entre otras condiciones, que ningún ensayo reservado
participe en el entrenamiento, que provenga de sujetos de entrenamiento, que
esté estratificado por clase y sesión, y que los estadísticos de normalización
se ajusten exclusivamente sobre el conjunto de entrenamiento reducido.

\subsection{Cota del efecto del cambio de régimen}
\label{sec:cota-confusor}

Para acotar cuánto del cambio entre la semilla $2026$ y las semillas $2027$ y $2028$ puede
deberse al régimen y no al azar, se calcula para cada sujeto $d_{\mathrm{cruce}} =
\mathrm{exac}(2026) - \tfrac12\bigl(\mathrm{exac}(2027) + \mathrm{exac}(2028)\bigr)$, con la
exactitud media del sujeto en cada semilla, y $d_{\mathrm{intra}} = \mathrm{exac}(2027) -
\mathrm{exac}(2028)$. Si el régimen no importara y el ruido de semilla tuviera la misma
varianza en las tres, ambas diferencias tendrían media cero, y la varianza de
$d_{\mathrm{cruce}}$ sería $0.75$ veces la de $d_{\mathrm{intra}}$. EEGMamba, que no cambia de
recorte, actúa como control; en el protocolo intra-sujeto, \eegnet{} y el Conformer, que
nunca tuvieron recorte, son un control sin cambio de recorte ni de hardware entre
semillas.

La cota del efecto sobre la exactitud media por sujeto es el mayor valor absoluto de los
extremos del intervalo al $95$\,\% (percentiles de un remuestreo de sujetos, $B = 10\,000$) de
la media de $d_{\mathrm{cruce}}$, o de su diferencia con la de EEGMamba si esa es mayor. En
una media de tres semillas, de las cuales una se entrenó sin recorte, la contaminación
máxima es la tercera parte de la cota. El análisis supone aditividad (el cambio de régimen
suma un desplazamiento constante a la exactitud de cada sujeto) y ruido de semilla
independiente; el cambio de régimen es compuesto (recorte, hardware y código) y no puede
separarse; y el control de EEGMamba supone que su sensibilidad al hardware y al código es la
de las demás arquitecturas, lo cual no es verificable. Es un análisis exploratorio y sus
pruebas no se corrigen por comparaciones múltiples.

El hardware de ambos regímenes difiere (tabla~\ref{tab:hardware}); no se interpreta ninguna
diferencia como efecto del hardware.


% =====================================================================
%  RESULTADOS
% =====================================================================

\section{Conclusiones}
Se compararon cuatro arquitecturas con $62$ canales: \eegnet{} con $2\,898$ parámetros, Conformer con $710\,146$, LaBraM Base con $5\,818\,138$ y EEGMamba con $3\,376\,445$. Solo LaBraM partió de pesos preentrenados.

En \loso{}, cada arquitectura aportó $54$ particiones por semilla y $162$ en total con tres semillas. La equivalencia práctica se examinó con margen $\pm 0.03$ y la prueba bayesiana incluyó márgenes de sensibilidad de $0.01$ y $0.02$.

La evaluación adicional combinó costo, estabilidad entre semillas, perturbaciones, interpretabilidad y calibración; los análisis de perturbaciones e interpretabilidad quedaron limitados a una semilla y diez sujetos apartados.
